\subsection{Organization and notation}\label{ms:guide}\label{ms:notation}
Section~\ref{v9:core:section} gives attainment, contact and reflection.
The next two sections develop the defect and the conservation identity.
These are the shared analytic tools. The remaining argument has two routes.
The compact route in Sections~\ref{v9:entry:section}--\ref{short:count:section}
gives strict exclusion at $n\ge700$ from a sharp predecessor.
The extension supplies a finite induction and a different local comparison
that applies at a first strict failure below that threshold as well.

For the logic of the extension, read the one-step estimate in
Section~\ref{ext:smoothing} and the first-failure argument in
Section~\ref{ext:bridge} before the numerical tables. Their conclusion
requires a sharp base through $42$ and two local exclusions.
Section~\ref{ext:starting-cases} treats the first four indices;
Sections~\ref{ext:finite-counts}--\ref{ext:floor40} describe the finite
induction through $42$. Section~\ref{ext:local-engine} explains the two
local exclusions. The continuous estimates and full numerical interfaces
are in the appendices. The assembly criteria state the numerical premises
explicitly, so each can be matched to its reproduction record.

Two optimization problems are used. The finite induction maximizes the
ratio, with attained score $C=C_n$. The compact tail argument maximizes
the excess above the fixed coefficient $c=\CE$. Both give contact,
but their excesses and normalization differ; Corollary~\ref{ext:ratio-attainment}
records the precise relation. Neither route changes the marginal law
while acquiring its successive bounds.

The following table distinguishes the three scales used in the proof.
Gap-specific parameters are introduced after the existence of a central
gap has been established.
\begin{center}\small
\begin{tabular}{@{}P{.27\linewidth}P{.66\linewidth}@{}}
\toprule
Symbol & Meaning and scale \\
\midrule
$H$, $H/\sqrt n$ & Threshold for the original sum; its normal-standardized value. \\
$B_0=\ind_{\{X>0\}}$, $p,q$ & Sign class and its probabilities for the original marginal. \\
$m$, $\ell=m/(pq)$ & $m=\E X_+=\E X_-$; distance between conditional means. \\
$W$, $\theta=\E W^2$ & $W=X-\E(X\mid B_0)$; conditional noise in original units. \\
$U=W/\ell$, $s=\E U^2$, $\lambda=ns$ & Noise in count units; marginal and total noise variances. \\
$t=np+H\sqrt{pq+s}=k+u$ & Count threshold, integer anchor, and phase. \\
$a_0,a_1$, $d=a_0+a_1$ & Inward support displacements in count units. \\
$J$, $D$ & Weighted support-gap sum and defect of the original law. \\
\bottomrule
\end{tabular}
\end{center}
The exact normalization is proved in \eqref{ms:bernoulli-coordinates}.
It identifies the count-coordinate law with the original $X$ and leaves
$H$ as its raw sum threshold. The signed moments are
$\gamma=\E X^3$ and $\beta=\E[X|X|]$.
